\chapter{Introduction}
\label{ch:introduction}

Thermal management is an essential part of spacecraft design, as the temperature
of each component must remain within allowable limits throughout the
mission. 

Since both the space environment and the operating conditions of the
spacecraft evolve across different mission phases, steady-state cases cannot fully represent the actual thermal response. 
Transient simulations are therefore needed to examine the evolution of temperatures and heat flows.

The design of a spacecraft thermal control system relies on numerical analysis
tools to predict system behaviour under different mission and operating
conditions. These tools are used to represent the spacecraft through a thermal
model and to calculate variables such as nodal temperatures and heat flows.
Some of the best-known tools include \acrshort{thermica}, \acrlong{td} together with \acrshort{sinda-fluint}, 
\acrfull{sst} and \acrshort{esatan-tms}, which provides a specialised 
environment for the development and solution of spacecraft thermal models.

The amount of information generated by these simulations can become
considerable, particularly when transient analyses or multiple cases of the
same model are considered. In such situations, the results must be organised
and represented in a way that allows their evolution to be examined and the
effects of the differences between cases to be identified. Dedicated
post-processing tools can therefore support the interpretation of the
simulation results.

This Bachelor’s Thesis builds on an existing Python-based post-processing tool
initially intended for steady-state analysis. The work combines the adaptation
of the tool to transient and multi-case analysis with the development of a
representative spacecraft thermal model in \acrshort{esatan-tms}. Selected
thermal results are also used to perform a simplified electrical power
analysis and provide the inputs for a separate battery model. The specific
objectives of the project are presented in the following section.

\section{Objectives}
\label{sec:objectives}

The main objective of the project is to enable the existing application to
process, analyse, and compare several transient cases associated with the same
\acrshort{esatan-tms} thermal model within a common environment.

Reaching this objective requires the coordinated development of the thermal
model and the post-processing tool. The \acrshort{esatan-tms} model provides
a consistent set of transient results, while the application is adapted to
manage and compare them. 

The thermal results are also used to estimate the electrical power generated
by the solar panels and to calculate the electrical power balance. These
calculations provide the inputs for a separate simplified battery model. 

Accordingly, the specific objectives of the project are:

\begin{itemize}
    \item Develop a representative spacecraft thermal model and define
    transient cases under different environmental and operating conditions.

    \item Adapt the existing application to load and manage several transient
    result files while retaining a common model hierarchy.

    \item Implement methods for representing and comparing temperatures,
    thermal loads, heat flows and exchanged thermal energy.

    \item Generate graphical outputs and summary tables, and allow the
    corresponding analyses and configurations to be saved, restored and
    exported.

    \item Estimate solar-panel electrical generation, calculate the electrical
    power balance and use the resulting power profiles as inputs to a separate
    simplified battery model.

    \item Apply the complete analysis process to the selected transient cases and compare
    the thermal responses obtained under different environmental and operating
    conditions.
\end{itemize}